\section{Factorization from the OPE}
\label{sec:ope}

In this section we make use of the operator product expansion to derive the matching relation for the quasi-PDF, as well as the equivalent matching relations for the \spcorr and pseudo-PDF. For simplicity these three equivalent cases are presented in separate subsections.

\subsection{OPE and Factorization for the Spatial Correlator}

The OPE is a technique to expand nonlocal operators with separation $z^\mu$ in terms of local ones in the Euclidean limit of $z^2\to 0$.
It can be applied to both bare regulated operators as well as renormalized operators, and our focus will be on the latter. For the gauge-invariant Wilson operator $\tilde{O}_\Gamma(z)$, it was proven that it can be multiplicatively renormalized in  coordinate space as~\cite{Ji:2017oey,Ishikawa:2017faj}
\beq \label{eq:ren}
\tilde{O}_\Gamma(z,\mu) = Z_{\psi,z}\, e^{\delta m |z|} \tilde{O}_\Gamma(z,\epsilon)\,,
\eeq
where $\delta m$ captures the power divergence from the Wilson line self-energy, $Z_{\psi,z}$ only depends on the end points $z,0$ and renormalizes the logarithmic divergences. This multiplicative renormalization was also discussed earlier in Refs.~\cite{Ishikawa:2016znu,Chen:2016fxx,Constantinou:2017sej}.  For simplicity, in this section we take $\Gamma=\gamma^z$ for \eq{ren}.

In the $\overline{\rm MS}$ scheme, the power divergence vanishes, and using the OPE
the renormalized $\tilde{O}_\Gamma(z,\mu)$ can be expanded
in terms of local gauge-invariant operators as $z^2\to 0$ giving
\begin{align}\label{eq:ope-tq}
\tilde{O}_{\gamma^z}(z,\mu) = & \sum_{n=0}^\infty \left[C_n ({\mu}^2 z^2)\frac{(-iz)^n}{n!} e_{\mu_1}\cdots 				e_{\mu_n}O_1^{\mu_0\mu_1\cdots\mu_n}(\mu)\right.
\\
&+C'_n ({\mu}^2 z^2)\frac{(-iz)^n}{n!} e_{\mu_1}\cdots e_{\mu_n}O_2^{\mu_0\mu_1\cdots\mu_n}(\mu)\nn\\
& +  \text{higher-twist operators}\Big]\,,\nn
\end{align}
where $\mu_0=z$, $C_n=1+O(\alpha_s)$ and $C'_n=O(\alpha_s)$ are Wilson coefficients, and $O_1^{\mu_0\mu_1\cdots\mu_n}(\mu)$ and $O_2^{\mu_0\mu_1\cdots\mu_n}(\mu)$ are the only allowed renormalized traceless symmetric twist-2 quark and gluon operators at leading power in the OPE,
\begin{align}\label{eq:twist-2}
O_1^{{\mu}_0 {\mu}_1 \ldots {\mu}_n}(\mu) = & Z_{n+1}^{qq} \bigl(
 \bar{\psi} \gamma^{({\mu}_0 } iD^{{\mu}_1} \cdots iD^{ {\mu}_n)} \psi -
\text{trace} \bigr) \,,
  \\
O_2^{{\mu}_0 {\mu}_1 \ldots {\mu}_n}(\mu) = & Z_{n+1}^{qg}\bigl(
 F^{(\mu_0\rho} iD^{{\mu}_1} \cdots iD^{ {\mu}_{n-1}} F_\rho^{~\mu_n)} - \text{trace} \bigr)
  .\nn
\end{align}
Here $Z_{n+1}^{ij}=Z_{n+1}^{ij}(\mu,\epsilon)$ are multiplicative $\overline{\rm MS}$ renormalization factors and $(\mu_0 \cdots \mu_n)$ stands for the symmetrization of these Lorentz indices.

The above OPE is valid for the operator itself, where we implicitly  constrain ourselves to the subspace of matrix elements for which the twist expansion is appropriate. In the iso-vector case, the mixing with the gluon operators is absent, which we will stick to for the rest of the paper. When $O_1^{\mu_0\mu_1\cdots\mu_n}$ is evaluated in the nucleon state,
\beq \label{eq:def-moments}
\langle P | O_1^{{\mu}_0 {\mu}_1 \cdots {\mu}_n} | P\rangle =  2a_{n + 1}(\mu)\left(P^{{\mu}_0} P^{{\mu}_1} \ldots P^{{\mu}_n} - \text{trace}\right), 
\eeq
where $a_{n+1}(\mu)$ is the $(n+1)$-th moment of the PDF,
\beq \label{eq:moment}
a_{n + 1} \left(\mu\right)=\int_{-1}^1 dx\,x^n q \left(x,\mu\right)\,,
\eeq
and the explicit expression of the trace term have been derived in Ref.~\cite{Nachtmann:1973mr,Georgi:1976ve,Chen:2016utp}. The inverse relation to \eq{moment} is that $q(x,\mu)$ has an expansion with terms proportional to the n'th derivative of the $\delta$-function, as in $\delta^{(n)}(x)\, a_n(\mu)$,  without any nontrivial short distance Wilson coefficient.

As pointed out in Ref.~\cite{Ji:2017rah}, to obtain enough information for the spatial correlator at $|\zeta|= |P^z z|\sim 1$ at small $z^2$, we have to choose $P^z$ large compared to the scale $\Lambda_{\rm QCD}$. When $P_z^2\gg\{\Lambda_{\rm QCD}^2, M^2\}$, the trace terms in Eq.(\ref{eq:def-moments}) are suppressed by powers of $M^2/P_z^2$, while the contributions from higher-twist operators in Eq.~(\ref{eq:ope-tq}) are suppressed by powers of $\Lambda_{\rm QCD}^2/P_z^2$ or $z^2\Lambda_{\rm QCD}^2$. Therefore, the twist-2 contribution is the leading approximation of the nucleon matrix element $\langle P|\tilde{O}_{\gamma^z}(z)|P\rangle$ at large momentum. From now on we will drop all the power corrections from our discussion.


The Wilson coefficients $C_n ({\mu}^2 z^2)$ in the OPE of $\tilde{O}_{\gamma^z}(z)$ can be computed in perturbation theory for $\mu\sim 1/|z|\gg \Lambda_{\rm QCD}$. In
the $\overline{\ensuremath{\operatorname{MS}}}$ scheme, the $C_n$ are log-singular near $z^2 =
0$, and so is $\langle P|\tilde{O}_{\gamma^z}(z,\mu)|P\rangle$. For this reason the $x$-moments of the quasi-PDF $\tilde q(x,P^z,\mu)$ are proportional to $C_n|_{z=0}$ which is
divergent, and the quasi-PDF will not simply become the PDF in the infinite $P^z$ limit. Instead, we need a factorization formula which matches the quasi-PDF to the PDF. In contrast, for the pseudo-PDF the moments do exist since we hold $\mu^2 z^2$ fixed when taking the $x$-moment. However, we still need a factorization formula to match the pseudo-PDF to the PDF.  We will comment further about this below.

Based on Eqs.~(\ref{eq:ope-tq}-\ref{eq:moment}), we can write down the leading-twist approximation to the spatial correlator as
\begin{align} \label{eq:fac-tQ-orig}
\tilde{Q}_{\gamma^z} \bigl(\zeta, {\mu}^2 z^2\bigr)
=&  \sum_{n} C_n ({\mu}^2 z^2)  \frac{(-i \zeta)^n}{n!} a_{n + 1}
 \left(\mu\right)
 \\
= & \sum_{n} C_n (\mu^2z^2)
\frac{(-i \zeta)^n}{n!}  \int_{-1}^1 dy\,y^n q \left( y, \mu\right)
 \,.  \nn
\end{align}
It should be noted that the only approximation we have made so far is ignoring
the higher-twist effects that are suppressed by small $z^2$ and the large momentum $P^z$ of the nucleon. In the limit of $P_z^2\gg M^2,\Lambda_{\rm QCD}^2$, we have $P^0\sim P^z$, so even if $\mu_0=0$ in Eq.~(\ref{eq:ope-tq}), the leading approximation of $\tilde{O}_{\gamma^0}(z)$ is still given by the twist-2 contributions in Eq.~(\ref{eq:fac-tQ-orig}), just with modified coefficients $C_n$.

Based on the OPE results in Eq.~(\ref{eq:fac-tQ-orig}), we can  derive a  factorization formula for the Euclidean \spcorr. First of all, let us
define a function $\mathcal{C} (\alpha, {\mu}^2 z^2)$:
\begin{align}  \label{eq:fac-C}
\mathcal{C} (\alpha, {\mu}^2 z^2) \equiv \int\! \frac{d \zeta}{2 \pi} \,
 e^{i \alpha \zeta}  \sum_n C_n ({\mu}^2 z^2)  \frac{(-i \zeta)^n}{n!}
 \,.
\end{align}
From \eq{fac-tQ-orig} and the renormalized analog of the Fourier-transform relation in \eq{pseudo} $\mathcal{C}(\alpha,\mu^2z^2)$ corresponds to a pseudo-PDF in the special case where $a_{n+1}(\mu)=1$. The analysis of Refs.~\cite{Radyushkin:1983wh,Radyushkin:2016hsy} implies that the support of $\mathcal{C}(\alpha,\mu^2z^2)$ is $-1 \le \alpha \le 1$. Noting that
\begin{align}
  \int d\alpha\, e^{-i\alpha(y \zeta)}\, \mathcal{C} (\alpha, {\mu}^2 z^2)
   &= \sum_n C_n(\mu^2 z^2) \frac{(-i\zeta y)^n}{n!} \,,
\end{align}
we find from \eq{fac-tQ-orig} that
\begin{align}
\tilde{Q}(\zeta,\mu^2 z^2)
= &\int_{-1}^1\!\! dy\int_{-1}^1\!\! d\alpha\ e^{-i\alpha(y\zeta)}\, \mathcal{C}(\alpha,\mu^2z^2)\, q(y,\mu) \,.
\end{align}
Finally, using the inverse transform of the renormalized analog of \eq{Qdefn},
\begin{align} \label{eq:FTq}
Q(\zeta,\mu) &= \int_{-1}^1\!\! dy\: e^{-i y\zeta}\,
  q(y,\mu) \,.
\end{align}
we obtain
\begin{align} \label{eq:io-Q-fact}
\tilde{Q}(\zeta,\mu^2 z^2)
=& \int_{-1}^1 d\alpha\ \mathcal{C}(\alpha,\mu^2z^2)\, Q(\alpha \zeta,\mu)
 + {\cal O}(z^2\Lambda_{\rm QCD}^2)
\,.
\end{align}
The result in \eq{io-Q-fact} is the factorization formula for the lattice calculable \spcorr $\tilde Q(\zeta,\mu^2 z^2)$ which expresses it in terms of the light-cone correlation $Q(\zeta,\mu)$ that defines the PDF. It has the same structure as the factorization formula for the \spcorr used for the calculation of the pion distribution amplitude in Ref.~\cite{Braun:2007wv,Bali:2017gfr}.

\subsection{Factorization for the quasi-PDF}

The renormalized quasi-PDF is defined as a Fourier transform of the renormalized \spcorr,
\begin{align} \label{eq:fac-tq-orig0}
\tilde{q} \Bigl(  x, \frac{\mu}{P_z}\Bigr)
\equiv &\int \frac{d \zeta}{2 \pi}\: e^{ix \zeta}\:  \tilde{Q} \biggl(\zeta, \frac{{\mu}^2\zeta^2}{P_z^2}\biggr) \,.
\end{align}
Note that we could use either $\tilde{Q}_{\gamma^z}$ or $\tilde{Q}_{\gamma^0}$ here.
Using the result for the \spcorr in Eq.~(\ref{eq:fac-tQ-orig}) this gives
\begin{align} \label{eq:fac-tq-orig}
\tilde{q} \Bigl( & x, \frac{\mu}{P_z}\Bigr)
 \\
= &  \int_{-1}^1 dy \biggl[ \int \frac{d \zeta}{2 \pi} e^{ix \zeta}
\sum_{n=0} C_n \Bigl( \frac{{\mu}^2\zeta^2}{P_z^2}\Bigr)  \frac{(-i
	\zeta)^n}{n!} y^n \biggr] q \left( y, \mu \right)
\nn\\
= &  \int_{-1}^1 {dy\over |y|} \biggl[\int \frac{d \zeta}{2 \pi} e^{i
	\frac{x}{y} \zeta}  \sum_{n=0} C_n \Bigl( \frac{{\mu}^2\zeta^2 }{(yP^z)^2}\Bigr)\,  \frac{(-i \zeta)^n}{n!}  \biggr] q \left( y,\mu \right) \,.
  \nn
\end{align}
Already, one can see that the matching kernel is a function of $x / y$
and $\mu / (|y|P^z)$. We define the kernel as
\begin{align}   \label{eq:fac-c}
C \left( {x\over y}, \frac{\mu}{|y|P^z} \right)
   \equiv \int\! \frac{d \zeta}{2 \pi} \: e^{i {x\over y} \zeta}
   \sum_{n=0} C_n \Bigl(
\frac{{\mu}^2\zeta^2}{(yP^z)^2}\Bigr) \, \frac{(-i \zeta)^n}{n!}
 \,,
\end{align}
and then Eq.~(\ref{eq:fac-tq-orig}) can be rewritten as
\begin{align} \label{eq:fac-tq}
\tilde{q} \left( x, \frac{\mu}{P^z}\right)
  = \int_{-1}^1 \frac{dy}{|y|}\: C \Bigl(
   \frac{x}{y}, \frac{\mu}{|y|P^z} \Bigr)\: q \left( y,\mu\right)
   \,,
\end{align}
which is the $\overline{\rm MS}$ factorization formula for the quasi-PDF. This result shows that the factorization formula in Eq.~(\ref{eq:momfact}) must have $p^z=|y|P^z$ for the quasi-PDF in the $\overline{\rm MS}$ scheme. We will show that this remains true for any quasi-PDF renormalization scheme in \sec{ren}. This differs
from the choice $p^z=P^z$ which had been conjectured and used in the early papers on the quasi-PDF~\cite{Ji:2013dva,Ji:2014gla,Xiong:2013bka}. Physically the correct result in \eq{fac-tq} can be understood as the fact that the matching coefficient is only sensitive to the perturbative partonic dynamics, and hence it is the magnitude of the partonic momentum $|y| P^z$ which appears, rather than the hadronic momentum $P^z$.

Taking the moment of the quasi-PDF using \eq{fac-tq-orig0} gives
\begin{align} \label{eq:momentqpdf}
 &\int_0^1\!\! dx\: x^n \tilde{q} \left( x, \frac{\mu}{P^z}\right)
= \Big(i\frac{d}{d\zeta}\Big)^n \tilde{Q} \biggl(\zeta, \frac{{\mu}^2\zeta^2}{P_z^2}\biggr) \Bigg|_{\zeta\to 0}
 \\
& \quad =  \sum_{n'}  \Big(i\frac{d}{d\zeta}\Big)^n \bigg[ C_{n'} \Big(\frac{{\mu}^2 \zeta^2}{P_z^2}\Big)  \frac{(-i \zeta)^{n'}}{n'!} \bigg] \Bigg|_{\zeta\to 0}  a_{n' + 1}
\left(\mu\right) \,.  \nn
\end{align}
Since the $C_{n'}$ coefficients have $\ln(\zeta^2)$ dependence, the derivative for $n'=n$ will always have a logarithmic singularity as $\zeta\to 0$, and there will be even more singular terms for $n'<n$. This explains why the short distance Wilson coefficient causes the moments not to exist for the quasi-PDF.


\subsection{Factorization for the pseudo-PDF}

The renormalized pseudo-PDF is the Fourier transform of the renormalized \spcorr
\begin{align} \label{eq:pseudoRen}
\mathcal{P} \left( x, \mu^2 z^2\right)
   = \int_{-\infty}^\infty \! \frac{d \zeta}{2\pi}\:
    e^{ix \zeta}\: \tilde{Q}\bigl( \zeta, \mu^2 z^2 \bigr)
   \,.
\end{align}
Since both the pseudo-PDF and \spcorr are multiplicatively renormalized in a $\zeta$-independent manner, this follows immediately from \eq{pseudo}.
If we take \eq{io-Q-fact} and Fourier transform the \spcorr $\tilde Q(\zeta,\mu^2z^2)$ into the pseudo-PDF, and light-cone correlation $Q(\alpha\zeta,\mu)$ into the PDF, then we immediately obtain the factorization formula for the pseudo-PDF,
\begin{align}\label{eq:ps-q-fact}
\mathcal{P}(x,z^2\mu^2) =
 &\int_{|x|}^1 {dy\over |y|}\
  \mathcal{C}\left({x\over y},\mu^2 z^2\right) q(y,\mu)\\
&+\int^{-|x|}_{-1} {dy\over |y|}\
  \mathcal{C}\left( \frac{x}{y}, \mu^2z^2\right) q(y,\mu)
  \nn \\
 &+ {\cal O}(z^2\Lambda_{\rm QCD}^2)\,,\nn
\end{align}
which is the small $z^2$ factorization formula in Eq.~(\ref{eq:pseudofact}). The upper and lower limits of the integrals in \eq{ps-q-fact} follow immediately from the support $-1\le \alpha \le 1$ of the matching coefficient $\mathcal{C}(\alpha,z^2\mu^2)$, and we recall that we also have $-1\le x\le 1$ for the pseudo-PDF on the LHS.

Since the range of $x$ is bounded for the pseudo-PDF the terms in the series expansion of the exponential in \eq{pseudoRen} exist,
\begin{align}
  \tilde Q(\zeta,\mu^2 z^2)
   = \sum_{n=0}^\infty \int_{-1}^{1}\!\! dx\: \frac{(-i\zeta)^n x^n}{n!}
     {\cal P}(x,\mu^2 z^2) \,.
\end{align}
Comparing with \eq{fac-tQ-orig} this implies that the moments of the pseudo-PDF are given by
\begin{align}  \label{eq:Pmoment}
  \int_{-1}^{1}\!\! dx\: x^n \:  {\cal P}(x,\mu^2 z^2)
  &= C_n(\mu^2 z^2)\, a_{n+1}(\mu) \,.
\end{align}

So far we have proven the large $P^z$ factorization of the quasi-PDF and small $z^2$ factorization of the spatial correlation and pseudo-PDFs. After
deriving one factorization, it immediately leads to the others, since they are just different representations of the same spatial correlator. Indeed, we see that the quasi-PDF and pseudo-PDF are related at leading power by their definitions:
\beq \label{eq:equiv}
\tilde{q}\Bigl(x,{\mu\over P^z}\Bigr)
  = \int\! {d\zeta\over 2\pi}\ e^{ix\zeta}\int_{-1}^1\!\! dy \ e^{-iy\zeta}
  \: \mathcal{P}\biggl(y,{\mu^2 \zeta^2\over P_z^2}\biggr)\,,
\eeq
where we have used $z=\zeta/P^z$.
Based on Eq.~(\ref{eq:fac-C}) and Eq.~(\ref{eq:fac-c}), the Wilson coefficients in their factorization theorems also maintain the same relationship,
\begin{align} \label{eq:ftc-C}
 C\left(\xi,{\mu\over |y|P^z}\right)
  &= \int\! {d\zeta\over 2\pi}\ e^{i\xi\zeta}\int_{-1}^1\!\! d\alpha\: e^{-i\alpha\zeta}\: \mathcal{C}\left(\alpha,{\mu^2 \zeta^2\over (yP^z)^2}\right) .
\end{align}
For the relations in \eqs{equiv}{ftc-C} the same choice of $\Gamma=\gamma^0$ or $\gamma^z$ should be used in the quasi- and pseudo-PDFs, or their corresponding coefficients.

In summary, there is a unique factorization formula that matches the quasi-PDF, \spcorr and pseudo-PDF to the PDF. Since their factorizations into the PDF all require small distances and have large nucleon momentum, the setup for their lattice calculations must also be the same. Therefore, the LaMET and pseudo distribution approaches are in principle equivalent to each other, and they differ perhaps only by effects related to their implementation on the lattice.

In Ref.~\cite{Radyushkin:2017cyf} it was speculated that one can study a ratio function
\beq \label{eq:ratio}
\tilde{Q}(\zeta,z^2,a^{-1})/\tilde{Q}(0,z^2,a^{-1})
\eeq
on a lattice with spacing $a$, and the $O(z^2)$ corrections may cancel approximately. This idea was tested in Ref.~\cite{Orginos:2017kos} in lattice QCD, and the results show that the ratio evolves slowly in $z^2$ at small values. It is then interesting to consider what type of non-perturbative information can be extracted from this ratio.

This question can be answered using the small $z^2$ factorization for the \spcorr.
 According to Eq.~(\ref{eq:fac-tQ-orig}),
\beq
	\tilde{Q}(0,\mu^2 z^2) = C_0(\mu^2z^2) + O(z^2\Lambda_{\rm QCD}^2)\,,
\eeq
where in the $\overline{\rm MS}$ scheme to one-loop
\beq  \label{eq:C0}
	C_0(\mu^2z^2) = 1 + {\alpha_sC_F\over 2\pi} \biggl[ {3\over2} \ln(\mu^2z^2 e^{2\gamma_E}/4)+{5\over2} \biggr] \,,
\eeq
which was also derived recently in \cite{Zhang:2018ggy}.
Then the ratio becomes
\begin{align} \label{eq:Qtratio}
	\frac{\tilde{Q}\left(\zeta, {\mu}^2 z^2\right)}{\tilde{Q}\left(0, {\mu}^2 z^2\right)}
	=&\sum_n \frac{C_n (\mu^2z^2)}{C_0 (\mu^2z^2)}	\frac{(-i \zeta)^n}{n!} a_{n+1}(\mu)\\
	&+ O(z^2\Lambda_{\rm QCD}^2)\,.\nn
\end{align}
Using \eq{Pmoment} and our $\overline{\rm MS}$ one-loop perturbative pseudo-PDF result in \eq{1looppseudopdfb} below we find for $\Gamma=\gamma^0$ that
\begin{align}
  C_n(\mu^2 z^2) &= 1 +\frac{\alpha_s C_F}{2\pi} \biggl[ \Bigl(
  \frac{3\!+\!2n}{2\!+\!3n\!+\!n^2}
  \!+\! 2H_n\Bigr) \ln\frac{\mu^2z^2 e^{2\gamma_E}}{4}
  \nn\\
 &\quad\quad
  +   \frac{5\!+\!2n}{2\!+\!3n\!+\!n^2}+2(1-H_n)H_n - 2 H_n^{(2)}
  \biggr] ,
\end{align}
where the Harmonic numbers are $H_n=\sum_{i=1}^n 1/i$ and $H_n^{(2)} =\sum_{i=1}^n 1/i^2$.  For the case $\Gamma=\gamma^z$ we have $C_n^{\gamma^z}(\mu^2 z^2) = C_n(\mu^2 z^2) + \Delta C_n^{\gamma^z}(\mu^2 z^2)$ with
\begin{align}
  \Delta C_n^{\gamma^z}(\mu^2 z^2)
   &= \frac{\alpha_s C_F}{2\pi} \: \frac{2}{2\!+\!3n\!+\!n^2} \,,
\end{align}
which also modifies \eq{C0} for $n=0$.
At small $z^2$ where the perturbative expansion with $\mu\simeq 1/|z|$ is valid, the ratio in \eq{Qtratio} has a weak logarithmic dependence on $|z|$, which is consistent with the lattice findings in Refs.~\cite{Orginos:2017kos,Karpie:2017bzm}. The weak dependence on $|z|$ can be quantitatively described by an evolution equation in $\ln z^2$~\cite{Radyushkin:2017cyf,Radyushkin:2017lvu}. According to our OPE analysis, here $\tilde{Q}(0,\mu^2 z^2)$ only serves as an overall normalization factor which is contaminated by higher-twist corrections, and the $\ln z^2$ evolution can be put in accurate terms with the factorization formula in Eq.~(\ref{eq:io-Q-fact}) that enables us extract the PDF from the ratio function. The same point was demonstrated by work done very recently in~\cite{Zhang:2018ggy}, which appeared simultaneously with our paper.
